\section{Literature Review}
\label{sec:lit}

\subsection{Analog synaptic devices}

An analog synapse needs a physical quantity that can be set to many stable values, read
without disturbing it, and moved between values at low energy. Five device families are
built around that requirement. They differ in how the state is stored, in what caps the
number of levels, and in how they fail.

Filamentary and non-filamentary oxide memristors (RRAM) set the weight by ion migration
that modulates an interfacial barrier or an oxygen-vacancy filament, and the two-terminal
structure this needs is simple enough to pack into dense crossbars: a
\ce{TiN}/\ce{TiOx}/\ce{HfO2}/\ce{TiN} array with $32\times36$ cells resolves eight stable
and distinguishable resistance states \cite{Hu2025_Sensors}. The same simplicity is what
the family pays for. The filament forms and ruptures by a stochastic nucleation
process each cycle, so cycle-to-cycle read noise and forming-voltage variability are
reported across the RRAM literature as a standing problem rather than an occasional one,
and a device-to-device separability statistic tied to a stated criterion, rather than a
level count simply asserted for a given pulsing protocol, is correspondingly hard to
find.

Phase-change memory (PCM) stores the weight as the amorphous-to-crystalline volume ratio
in a chalcogenide alloy, and of the five families it has the most mature manufacturing
record: a 128~Mb \ce{Ge2Sb2Te5} chip has been fabricated at a memory yield of
\SI{99.99999}{\percent} \cite{Xie2025_AdvSci}, and a 64-core mixed-signal in-memory
compute chip built on the same material reaches \num{63.1}~TOPS of throughput at
\num{9.76}~TOPS/W with accuracy comparable to a software implementation on convolutional
and LSTM layers \cite{LeGallo2023_NatElectron}. Programming energy has been pushed low: a
graphene-nanoribbon edge contact with a cross-sectional area of about \SI{1}{\nano\metre
\squared} cuts the write power to about \SI{53.7}{\femto\joule} per pulse
\cite{Wang2022_AdvSci_PCM}. The known failure mode is conductance drift, which follows a
power law $G(t)=G(t_0)(t/t_0)^{-\nu}$ with a phase- and material-dependent exponent
$\nu$, and a second, structural asymmetry between a gradual SET and an abrupt RESET drives
weight-update asymmetry severe enough that in-situ trained accuracy has generally trailed
software training, which is why differential $G^+-G^-$ synapse pairs are the standard
mitigation \cite{Ambrogio2018_Nature}. PCM is also the exception to the separability
problem that runs through this survey. Joshi et al.\ demonstrate 11 iteratively
programmed conductance levels and report the statistic behind that count directly: a
standard deviation below \SI{1.2}{\micro\siemens} for every level, measured across
\num{10000} devices per level \cite{Joshi2020_NatComm}. That is a measured population
statistic tied to an explicit criterion, not an asserted number, and it is close to
singular in this literature for being so. The same devices, mapped into a differential
configuration at two PCM cells per weight, reach \SI{93.7}{\percent} on CIFAR-10 and
\SI{71.6}{\percent} top-1 on ImageNet after training-weight transfer \cite{Joshi2020_NatComm}.

Electrochemical RAM (ECRAM) sets the channel conductance non-volatilely by
electrochemically controlled intercalation of a mobile ionic species (protons, in the
device that gives this family its best-verified numbers) into the channel material,
under fields of order \SI{1}{\volt\per\nano\metre} \cite{Onen2022_Science}. Because the
write and read paths are electrically separate, potentiation and depression are reported
as close to linear and symmetric, with low-voltage switching and low power dissipation
\cite{Fuller2017_AdvMat}. ECRAM's strongest reported figure is energy: proton transfer during a
\SI{5}{\nano\second} write pulse costs about \SI{15}{\atto\joule}
\cite{Onen2022_Science}, roughly three orders of magnitude below PCM's best verified
figure, and a separate CMOS-compatible device confirms the same \SI{5}{\nano\second}
switching speed while reaching \SI{96}{\percent} accuracy on an MNIST simulation built
from its own measured device characteristics, with a projected switching energy near
\SI{1}{\femto\joule} at a larger $100\times100$~nm$^2$ device geometry
\cite{Tang2018_IEDM}. Endurance is also strong: an ionic floating-gate array endures
more than one billion write-read operations at write-read frequencies above
\SI{1}{\mega\hertz} \cite{Fuller2019_Science}. What ECRAM has not produced, across two
dedicated literature passes for this review, is a verified separability criterion behind
any of its level-count claims. The strongest verified numbers are a dynamic-range ratio
of $20\times$ with no discrete count attached \cite{Onen2022_Science} and a device that
demonstrates ``up to 1000 discrete states'' \cite{Tang2018_IEDM} with no accompanying
$\sigma/\mu$ or read-margin statistic, an asserted number in the same sense as most of
the RRAM literature above.

A ferroelectric FET stores its state as remanent polarization, and because the
switching mechanism is domain nucleation and growth rather than ion transport it is
reported at low programming energy: an HZO-gated device coupled to a two-dimensional
\ce{ReS2} channel switches at \SI{2.0}{\femto\joule} per event
\cite{Huo2026_NatComm}. The known failure modes are interfacial: instabilities and
depolarization fields at the two-dimensional-channel/ferroelectric junction degrade both
the retained polarization and long-term retention in that same device
\cite{Huo2026_NatComm}, and an ANN and CNN built from it reach \SI{97.33}{\percent} on
iris recognition and successful polarization-resolved butterfly classification
respectively \cite{Huo2026_NatComm}. Level counts in the FeFET literature are taken up in
Section~\ref{sec:lit:gaa} below, where they belong with the programming schemes that
produce them; as with the families above, they are stated for a given pulsing protocol
rather than against a measured separation criterion.

Ferroelectric tunnel junctions (FTJ) store the weight in tunnelling electroresistance
rather than channel conductance: polarization reversal changes the effective tunnelling
barrier at a metal/ferroelectric/metal or metal/ferroelectric/semiconductor junction, and
the on/off ratio this produces is reported rather than a discrete level count.
Reliability issues are the same interface story as the FeFET case: retention instability
from the depolarization field, and endurance weakened by the field concentrated at the
ferroelectric/dielectric interface during program/erase cycling
\cite{You2026_AdvSci}. Where FTJ synapses have been carried through to a network
demonstration, forming-gas annealing that passivates interfacial defects gives
sufficiently linear and symmetric potentiation and depression that a pattern-recognition
simulation on a \ce{HfOx} crossbar reaches about \SI{90}{\percent} accuracy
\cite{Nguyen2025_NanoConverg}.

Read across all five families, the pattern is consistent: a level count is usually stated
for a fixed pulsing protocol, and a measured spread tied to an explicit separability
criterion is the exception rather than the rule. Joshi et al.'s PCM result (eleven
levels, a standard deviation under \SI{1.2}{\micro\siemens}, measured across ten thousand
devices per level \cite{Joshi2020_NatComm}) is worth naming as the exception precisely
because it is rare enough to make the general point credible rather than convenient: most
of this literature reports a number, not the statistic that would justify calling it a
level. Section~\ref{sec:device} follows Joshi et al.'s example rather than the majority
pattern, reporting a conservative eight-level count under its own explicit three-sigma
separation criterion.

\subsection{Ferroelectric \ce{HfO2} and HZO}

Ferroelectricity in doped \ce{HfO2} was reported by B\"oscke et al.\ in 2011: a
\SI{10}{\nano\metre} film with less than \SI{4}{\percent} \ce{SiO2} doping crystallizes
under electrode encapsulation into a phase mixture with a remanent polarization above
\SI{10}{\micro\coulomb\per\square\centi\metre} at a coercive field of
\SI{1}{\mega\volt\per\centi\metre} \cite{Boescke2011_APL}. That the effect appears in
\ce{HfO2} rather than a perovskite mattered for reasons beyond novelty. Perovskite
ferroelectrics such as \ce{PbZrTiO3} lose polarization as film thickness scales down and
are, in their own right, incompatible with complementary-metal-oxide-semiconductor
processing \cite{Liao2023_FundamentalResearch}; \ce{HfO2} is already a qualified
high-$\kappa$ gate dielectric, ferroelectric down to a few nanometres, and compatible with
a back-end-of-line thermal budget. A TiN-capped \ce{Hf0.5Zr0.5O2} (HZO) capacitor
crystallized by stress rather than by high-temperature annealing reaches a switching
polarization of \SI{45}{\micro\coulomb\per\square\centi\metre} at a ferroelectric
saturation voltage near \SI{1.5}{\volt} \cite{Kim2017_APL}.

The polar phase responsible is the orthorhombic $Pca2_1$ structure, distinguished from
the related non-polar $Pmn2_1$ orthorhombic phase by symmetry and by its response to
shear strain \cite{Chen2023_APL}. It is stabilised by dopant choice, capping-layer stress,
grain size and anneal condition, and interface engineering has pushed the polar phase down
to a film as thin as \SI{1.5}{\nano\metre} \cite{Shi2023_NatCommun}. In the
\SIrange{5}{10}{\nano\metre} range relevant to a gate stack, reported remanent
polarization runs from about \SI{10}{\micro\coulomb\per\square\centi\metre} in the
earliest films \cite{Boescke2011_APL} up to $2P_r\approx\SI{42}{\micro\coulomb\per\square
\centi\metre}$ in an interlayer-engineered capacitor \cite{Hsain2023_Nanotechnology}, at
coercive fields of order \SI{1}{\mega\volt\per\centi\metre}. At the thin end of that
range the film's own electrostatics start to compete with switching: in \SI{5}{\nano
\metre} HZO, depolarization fields drive a stable paraelectric phase to coexist with a
metastable ferroelectric one, which is what produces a pinched hysteresis loop rather than
a saturated one \cite{Siannas2022_CommPhys}.

None of this is fixed once a device is written. Wake-up, fatigue and imprint are the
three reliability effects that determine what a written state looks like after the write
pulse is over, and each has a distinct mechanism, magnitude and timescale.

Wake-up is the polarization increase seen over the first cycles of operation, and its
mechanism has been imaged directly: operando plasmon-enhanced spectroscopy, sensitive to
under \SI{1}{\percent} oxygen-vacancy variation, tracks vacancy migration in
\SI{5}{\nano\metre} HZO during the pre-wake-up stage in real time \cite{Jan2023_AFM}. It
is a transient effect, not a permanent one: devices wake up for up to $10^3$ bipolar
pulsed-field cycles \cite{Fields2021_JAP}. Beyond that cycle count the same mechanism
runs in reverse. A critical field of about \SI{0.8}{\mega\volt\per\centi\metre} separates
wake-up from fatigue by governing a reversible transition between polar and antipolar
phases \cite{Cheng2022_NatCommun}, and oxygen-vacancy concentration and redistribution
are the correlated driver of both switching speed and endurance more generally
\cite{Lee2023_NanoConvergence}. Fatigue is reported with a magnitude and a cycle count
together: it drives the polarization toward zero
by $10^8$ cycles \cite{Fields2021_JAP}, and an independent measurement on an
alumina-interlayer capacitor corroborates an endurance limit beyond $10^8$ cycles at a
cycling field of \SI{3.5}{\mega\volt\per\centi\metre}, with $2P_r$ held near
\SI{42}{\micro\coulomb\per\square\centi\metre} up to that point
\cite{Hsain2023_Nanotechnology}. Doping is a reported partial fix: (La,Y)- and
(La,Gd)-co-doped hafnium zirconate capacitors are fatigue-free, with remanent
polarization at or above \SI{15}{\micro\coulomb\per\square\centi\metre}, out to $10^{11}$
cycles \cite{Popovici2022_ACSAEM}.

Imprint is an asymmetric shift of the hysteresis loop along the voltage axis, typically
attributed to asymmetric charge trapping that builds an internal bias field. It is the
least well documented of the three effects in the verified literature. Its mechanism has
been isolated by
bake-time- and bake-temperature-dependent current-peak splitting in a La-doped HZO
capacitor, and the authors extract an activation energy consistent with oxygen-vacancy
redistribution driving what they term a wake-up reversal \cite{Lee2023_IEEETED}; the
magnitude that activation energy corresponds to, in volts of threshold shift or percent of
$P_r$, is not established by a source this review can verify, and no number is asserted
here for it.

Depolarization and charge trapping are a second pair of effects from the same
interface, and they overlap. In a metal-ferroelectric-insulator-semiconductor stack,
incomplete screening
of the ferroelectric's bound charge by the interfacial layer produces a depolarization
field that opposes the stored polarization; combined with charge trapping at the
ferroelectric/interfacial-oxide interface, the two effects are strongly coupled and can
mask each other in a threshold-voltage or capacitance-voltage measurement
\cite{Toprasertpong2022_ApplPhysA}. Uncorrected, that combination sets the
retention floor: a device engineered to isolate the two contributions on the same
polarization-voltage and current-voltage measurement reaches a memory window of about
\SI{0.8}{\volt} at \SI{85}{\celsius}, extrapolated to ten years
\cite{Higashi2021_IEEETED}. The same depolarization physics that pins the memory window
in that stack produces the pinched, unsaturated loop in thin HZO noted above
\cite{Siannas2022_CommPhys}. One mechanism appears as two different symptoms,
depending on how thick the film is and what it is stacked against.

\subsection{FeFET synapses and the gate-all-around architecture}
\label{sec:lit:gaa}

Published FeFET analog synapses cluster around what identical-pulse and
linearity-engineered programming schemes can each deliver. Under domain
nucleation-limited switching driven by trains of identical pulses, a five-bit synapse
(32 conductance states across a $45\times$ tunable range) is reported at a
\SI{75}{\nano\second} update pulse \cite{Jerry2017_IEDM}. Redesigning the gate stack as a
ferroelectric/dielectric superlattice to linearize the weight-update profile pushes this
to seven bits, 128 states, with \SI{94.1}{\percent} accuracy on an online-trained MNIST
task and endurance beyond $10^{10}$ cycles \cite{Aabrar2022_TED}. At array scale, a
multi-level-cell FeFET crossbar built in a 28~nm high-$\kappa$/metal-gate process reaches
\num{885.4}~TOPS/W and \SI{96.6}{\percent} handwriting-recognition and
\SI{91.5}{\percent} image-classification accuracy \cite{Soliman2023_NatComm}. None of
these three reports states a separability criterion behind its level count in the sense
Section~\ref{sec:levels} below will use one; each states a count achieved under a
specific pulsing or verification protocol, which is a weaker claim.

Gate-all-around, nanosheet-channel geometry is the industry's response to the same
electrostatic-control problem an analog FeFET synapse also has to solve: a channel that
is fully surrounded by gate holds its threshold and subthreshold behaviour against drain
bias and channel-length variation better than one gated on three sides. A dual-$\kappa$
spacer multi-channel nanosheet device reports a subthreshold swing of
\SI{62}{\milli\volt\per\text{dec}}, drain-induced barrier lowering of
\SI{14}{\milli\volt\per\volt} and an on/off ratio near $10^5$
\cite{Panigrahy2024_IEEEAccess}. A benchmarking study against single-channel FinFETs
at equal effective width finds that a multi-bridge-channel nanosheet reduces propagation
delay by \SI{31}{\percent} over a nanowire predecessor and a further \SI{12}{\percent}
over that, with a \SI{56}{\percent} ring-oscillator frequency gain, and improves
on-current, on/off ratio, subthreshold swing and DIBL at the same time
\cite{Sreenivasulu2024_IEEEAccess}.

The pairing of a ferroelectric gate stack with a gate-all-around channel is already
published work, on both counts. Stacked two-nanosheet gate-all-around \ce{Ge0.98Si0.02}
FeFETs reach a \SI{1.8}{\volt} memory window at a \SI{2}{\volt} write voltage and
endurance beyond $10^{11}$ cycles at a \SI{400}{\celsius} thermal budget
\cite{Chen2023_VLSI}, an interfacial-layer-free follow-on of the same device family
reproduces those figures with a \ce{Ge0.95Si0.05} channel \cite{Hsieh2024_TED}, and a
three-dimensional double-HZO gate-all-around nanosheet device built specifically to
homogenize the corner field reports a memory window of \SIrange{0.9}{1.3}{\volt}
depending on interlayer choice, at a program/erase voltage of
\SI{\pm3.5}{\volt} and endurance beyond $10^{11}$ cycles \cite{Liao2022_VLSI}.
Gate-all-around nanowire cross-sections carry the same pairing: a stacked, sub-5-nm HZO
nanowire FeFET with an internal metal gate reaches a minimum subthreshold swing of
\SI{49.3}{\milli\volt\per\text{dec}} \cite{Lee2021_JEDS}, and a vertical III-V nanowire
gate-all-around channel on silicon has also been integrated with an HZO gate stack
\cite{Persson2022_EDL}. As reported, these six devices are characterized on memory-window,
retention, endurance and read-current metrics; none of their published abstracts states a
conductance-level count, a linearity or symmetry figure, or any other analog-synapse
characterization, so they should be read as ferroelectric gate-all-around \emph{memory}
demonstrations rather than as ferroelectric gate-all-around \emph{synapse}
demonstrations.

The one exception is also the closest prior work to the device studied in this thesis.
Lin and Su model a stacked-nanosheet FeFET explicitly as a synapse: Monte Carlo
simulation under a nucleation-limited-switching ferroelectric model, considering
cycle-to-cycle variation, of the intrinsic conductance response under identical-pulse
stimulation \cite{Lin2024_OJNano,Lin2023_SNW}. Increasing the number of stacked channel
tiers improves the maximum-to-minimum conductance ratio and the per-state noise margin
without a footprint penalty, and a companion study attributes the same tier-count benefit
to reduced variability in threshold voltage, memory window and synapse conductance
response, sourced from random ferroelectric-dielectric phase distribution rather than
switching stochasticity \cite{Lin2023_SNW}. What this pair of papers does not report,
so far as their published text states, is a discrete conductance-level count, a
programming energy, or a neural network trained and deployed on the resulting response.
A ferroelectric gate-all-around synapse is accordingly published work; a ferroelectric
gate-all-around synapse calibrated to a measured device, with a measured write-energy
budget, a measured cycle-to-cycle statistic and a deployed network, is not.

\subsection{Spiking networks for biosignal classification}

The architecture and task this thesis's own network adopts originate with Yuan et al.,
whose neuromorphic system processes physiological signals through a leaky-and-adaptive
spiking neural network (LSNN) with VO$_2$-memristor spike encoding. On a curated MIT-BIH
set of 2000 heartbeats (1000 normal, 500 ventricular-ectopic, 250 supraventricular-ectopic
and 250 fusion beats, split randomly into 1664 training and 336 test beats), a
$3\times100\times4$ LSNN with 60 leaky-integrate-and-fire (LIF) and 40 adaptive-LIF (ALIF)
hidden neurons reaches \SI{95.83}{\percent} on the four-class arrhythmia task and
\SI{99.79}{\percent} on a separate epileptic-seizure-detection task
\cite{Yuan2023_NatComm}. The mixed LIF/ALIF population is not incidental: an LIF-only
network of the same size underperforms the mixed network by about \SI{15}{\percent} on a
two-class version of the task and about \SI{18}{\percent} on the four-class version
\cite{Yuan2023_NatComm}, which is why the adaptive population is part of the architecture
this thesis carries forward rather than a detail dropped in translation.

Accuracy figures in this literature are not directly comparable across papers, and the
class count has to travel with each one. A neuromorphic classifier trained on a
different, much larger clinical resource (45{,}152 patients, twelve-lead recordings)
reaches \SI{94.4}{\percent} overall accuracy on a four-class task, but the four classes
are atrial fibrillation, sinus rhythm, and left- and right-bundle-branch block rather than
the AAMI normal/ventricular/supraventricular/fusion scheme, and the split is strict
inter-patient rather than the pooled random split used above \cite{Kolhar2025_SciRep}.
Despite both papers reporting ``four classes,'' the two accuracies are not measurements of
the same problem and should not be placed on the same axis. A two-class SNN on joint ECG
and phonocardiogram signals, evaluated on a patient-wise inter-patient split of the
PhysioNet/CinC 2016 challenge set, reaches \SI{89.74}{\percent} with a reported energy
consumption of \SI{209.6}{\micro\joule} per classification \cite{Ran2025_Sensors}. That
is a binary task, and correspondingly not comparable to either four-class number
above.

Deployment onto physical neuromorphic hardware is thin in the ECG-specific literature
this review located. The one Loihi energy figure surfaced in this search, \SI{23.6}{\pico
\joule} per spike, traces back to Intel's own published chip specification, cited by an
unrelated transparent-healthcare paper rather than measured on an ECG classification
workload \cite{Sungheetha2026_SciRep}; no source verified for this review reports an
ECG-specific spiking classifier measured on Loihi, SpiNNaker, TrueNorth or a comparable
neuromorphic chip.

On what a trained network loses when quantised onto a small number of analog conductance
levels, the verified literature offers mechanism rather than a number. Nonlinear or
asymmetric weight updates are reported to introduce gradient distortion and to degrade
learning accuracy in neuromorphic systems built on multi-level conductance devices
\cite{Kumar2026_Nanomaterials}, and device variability, retention drift and non-ideal
conductance updates are named as the limiting factors behind the still-small number of
demonstrated multi-level systems in two-dimensional-material devices
\cite{Kazmi2026_NanoMicroLett}. Neither source, nor any other reached by this review,
attaches a quantised-accuracy-loss percentage to a stated level count, and none addresses
whether that loss depends on how the levels are spaced rather than merely how many of
them there are. That question is left open by the literature surveyed here.

What this review leaves open, across all four subsections, is what Section~\ref{sec:method}
takes up next: a device calibrated against a measured transfer characteristic rather than
against its own simulation, a measured write-energy and cycle-to-cycle noise budget rather
than a modelled one, level placement examined as a quantity distinct from level count, and
a trained network actually deployed onto the resulting conductance ladder rather than onto
an idealized one.
